%==============================================================================
%  Chapter 6 — Conclusion
%
%  Short, decisive chapter (typically 4--7 pages):
%   - restate the question
%   - summarise findings
%   - state contributions
%   - acknowledge limitations
%   - propose future research
%   - close
%==============================================================================
\chapter{Conclusion}
\label{ch:conclusion}

\section{Restating the research question}
\label{sec:conc-rq}

The thesis set out to answer one main question and four sub-questions.
The main question is restated verbatim from \Cref{sec:intro-rq}:

\medskip
\noindent\textit{How are investment logics constructed, legitimised and
contested within the EIC instruments Pathfinder and STEP, and how can
these dynamics be understood through the lens of institutional theory?}
\medskip

The sub-questions framed the inquiry as a textual one. SQ1 asked which
logics the Pathfinder and STEP texts of the EIC Work Programme 2026
inscribe, and how those logics are distributed across each instrument's
architecture. SQ2 asked how the text constructs legitimacy for the
co-presence of logics that the theoretical literature treats as
divergent. SQ3 asked how STEP reconfigures the relationship between
state and market logic relative to Pathfinder. SQ4 asked what the
comparison reveals about the model of public investment the EU is
institutionalising. The empirical scope was documentary throughout: the
Pathfinder and STEP Scale Up sections of the Work Programme, read
against the regulatory instruments that establish them, with no
interview material.

\section{Summary of findings}
\label{sec:conc-findings}

\paragraph{SQ1: inscription and distribution of logics.}
The corpus inscribes four logics in Thornton's sense: science, market,
state, and profession. They are distributed by pipeline position, not
by what applicants do. Pathfinder is organised primarily by a science
logic at TRL~1--4, with state-defined Challenge relevance embedded as a
sub-criterion of scientific excellence; STEP Scale Up is organised
primarily by a state logic operating through market mechanics at the
scale-up end of the pipeline. The macro-level succession Science\,$\to$
Market\,$\to$\,State is reproduced inside every Challenge chapter as
Pattern~D, so that substantively different technology domains receive
the same four-step logic order. The reproduction is a feature of the
chapter template rather than a derivation from the science.

\paragraph{SQ2: legitimacy work for the co-presence of logics.}
The Work Programme legitimises co-presence through three distinct
strategies. Pattern~A (tension-dissolution) joins logics by syntactic
juxtaposition rather than by argument, recruiting the
taken-for-grantedness of bureaucratic prose itself; Pattern~F shows the
Technology Readiness Level framework carrying three legitimation axes
within a single designed artefact, so that the political question of
where public money should flow in the pipeline is presented as a
technical question of developmental readiness; the Sovereignty Seal
operates as a designed credential that carries state-issued legitimacy
across institutional fields independently of the capital decision that
produced it. The three strategies operate at different scales
(sentence, artefact, credential) and together perform the bulk of the
text's legitimation work.

\paragraph{SQ3: STEP and the state--market reconfiguration.}
STEP is best read not as the EU adopting market venture-capital logic
but as a developmental-state instrument deploying market mechanics.
Pattern~E (circular market validation) shows the simultaneous
requirements of demonstrated market interest and demonstrated market
failure as satisfiable only through the EIC's prospective anchor
participation: the state constructs the market signal it then uses as
legitimation for its own intervention. The configuration is consistent
with the Chinese Government Guidance Funds model and with
\textcite{mazzucato2018mission}'s patient-capital account; it is
inconsistent with the surface vocabulary of market gap and deal flow in
which the instrument narrates itself. The interpretive consequence
captured by the \emph{logic-as-register} distinction is that STEP is
best classified by its structure rather than its vocabulary.

\paragraph{SQ4: the EU model of public investment.}
The comparison reveals an instrument whose developmental-state
character is held in place by three further features beyond Pattern~E.
Pattern~B shows coercive isomorphism operating through grant-agreement
mechanics: the autonomy promise at [021] is real at the level of
methodological choice and sharply bounded at the level of direction by
the Programme Manager's contractually authorised override authority.
Pattern~C shows a state social-equity template inserted into the
tiebreaker architecture at [059], ranking gender and inclusion
priorities above the instrument's own market and geographic criteria.
A candidate fifth template for sustainability was tested in Pass~2 and
not confirmed: sustainability enters through four heterogeneous
mechanisms rather than through a universal insert. The EU model that
emerges is one in which mission-oriented investment is carried by an
architecture that templates cross-cutting normative priorities,
mediates between political principal and applicant through hired
technocrats with contractual discretion, and combines borrowed market
vocabulary with state-anchored structural design.

\section{Contributions revisited}
\label{sec:conc-contrib}

The contributions stated in \Cref{sec:intro-contribution} and developed
in \Cref{sec:disc-contrib} can now be restated against the evidence
that supports them.

\paragraph{Theoretical.}
Four contributions to the institutional-logics literature on
supranational public-investment instruments. The
\emph{logic-as-register} distinction (Pass~2 coding rule~7) separates
the vocabulary an instrument speaks from the logic that organises it,
and is generally useful wherever the state borrows market discourse
without transferring authority to market actors. \emph{Tension-dissolution}
(Pattern~A) names the textual move that presents conflicting logics as
naturally complementary, and merits naming inside Suchman's
cognitive-legitimacy typology as distinct from surface-visible logic
conflict and from \textcite{friedland1991bringing}'s mobilisation of
contradiction. The Sovereignty Seal is advanced as a portable
field-bridging legitimacy artefact: a designed credential whose
legitimacy is engineered to travel across institutional fields,
decoupled from the capital decision that produced it. \emph{Circular
market validation} (Pattern~E) names the configuration of simultaneous
market-interest and market-failure requirements satisfiable only
through state-anchor participation, and is offered as a structural
property of the class of developmental-state instruments to which STEP
belongs. A fifth, more modest observation is that Pattern~F's
triple-load shows \textcite{scott2014institutions}'s three pillars
co-operating within a single designed artefact, refining rather than
contradicting Scott's frame.

\paragraph{Empirical.}
The thesis offers the first comparative qualitative reading, to the
author's knowledge, of Pathfinder and STEP Scale Up as co-inscribed in
the EIC Work Programme 2026 that works in a critical institutional
register rather than in the policy-advocacy register dominant in the
wider literature on EU blended finance. The reading is anchored in a
numbered, traceable corpus of coded entries (180 Pass~1 entries, 29
Pass~2 resolutions, six confirmed patterns, one candidate not
confirmed), with the codebooks reproduced in
\Cref{app:codebook-deductive,app:codebook-inductive} and the audit
trail in \Cref{app:audit-trail}.

\paragraph{Practical.}
The reading identifies the specific textual mechanisms by which the
Work Programme manages logic conflict and locates them at points an
FP10 successor instrument can adjust deliberately: the tiebreaker
architecture, the Excellence sub-criterion embedding, the
grant-agreement override clauses, and the Seal's decoupling from the
investment decision. The recommendation throughout
\Cref{sec:disc-practice} is for explicitness about the trade-offs
already present in the text rather than for their resolution.

\section{Limitations}
\label{sec:conc-limits}

The methodological limitations of the study (single coder, single
corpus, reflexive position) are addressed in \Cref{ch:methodology}.
What follows are substantive limitations of the findings themselves.

The corpus is a single Work Programme year. The 2026 Work Programme is
the first to incorporate STEP Scale Up in its current form, and the
configuration it inscribes may shift as the instrument matures. The
findings should be read as a reading of the institutional design as it
stands in 2026, not as claims about how the design will stabilise.

The instrument scope excludes the EIC Accelerator and the EIC
Transition instrument. Both sit between Pathfinder and STEP Scale Up in
the pipeline and use distinct evaluation architectures. The claim that
the EIC distributes logics by pipeline position is supported across the
two endpoints studied; whether the pattern holds continuously across
the middle of the pipeline is left open.

The thesis reads only what the text inscribes. The reading captures
what the Work Programme commits the instruments to, not what Programme
Managers, evaluators, applicants, or investee firms then do with those
commitments. Pattern~B in particular is consequential at the level of
contractual authority structure, and its operational character cannot
be settled from the text alone.

The four non-EU comparators (DARPA, Temasek, the Chinese Government
Guidance Funds, and the Mazzucato patient-capital model) appear only
where the Work Programme invokes them or where the secondary literature
uses them to characterise mission-oriented design. They are referenced
rather than subjected to the same close reading as the EIC text, and
the mimetic-isomorphism reading at Pattern~F and at
\Cref{sec:disc-step-mimesis} relies on the secondary literature for the
characterisation of the referent fields.

The firm-level observation that motivated the inquiry
(\Cref{sec:intro-origin}) is not tested by the thesis. The textual
reading is consistent with a public-instrument architecture that could
produce the under-capitalisation dynamic the prior view describes, but
the dynamic itself is a firm-level claim that requires firm-level data
and a different research design.

\section{Avenues for future research}
\label{sec:conc-future}

The findings open four lines of further inquiry.
\begin{itemize}[leftmargin=*]
    \item \textbf{Quantitative follow-up on funded portfolios.} A
          firm-level study tracking Pathfinder and STEP awardees through
          subsequent private financing rounds would test whether the
          architectural features identified in the text correspond to
          the financing trajectories the origin observation describes.
          The textual reading provides the design hypotheses; the
          firm-level data would test them.
    \item \textbf{Comparative work with non-EU innovation agencies.}
          Close textual readings of ARIA in the United Kingdom, ARPA-E
          in the United States, and JST in Japan, applying the same
          coding architecture, would establish whether the
          tension-dissolution, circular market validation and portable
          credential moves are EU-specific or recur across coordinated
          mission-oriented designs. The Pattern~A scope condition
          (industrial-policy vocabulary being contested) is the most
          immediate object such a comparison can adjudicate.
    \item \textbf{Longitudinal study as STEP matures.} STEP Scale Up is
          new. A repeat reading of the 2028 and 2030 Work Programmes
          would test whether the developmental-state inflection
          consolidates, whether the Sovereignty Seal's portability is
          institutionalised across downstream EU vehicles, and whether
          the circular-validation configuration is retained or modified
          as private co-investors accumulate experience with the
          instrument.
    \item \textbf{Cross-instrument and cross-pillar extensions inside
          the EU.} A reading of the Innovation Fund, an EU Sovereignty
          Fund if established, and the EIC Accelerator and Transition
          instruments, using the same codebook, would test how widely
          the templates identified here travel inside the EU
          architecture. The Pattern~D logic succession and the
          Pattern~C tiebreaker template are the most portable
          candidates; the Pattern~F triple-load travels wherever TRL
          performs a regulative gate function.
\end{itemize}

\section{Closing remarks}
\label{sec:conc-close}

The Work Programme is not only a funding document. It is the textual
artefact through which the EU constitutes who counts as an EIC
applicant, what counts as a credible innovation claim, and what counts
as a successful outcome. The reading offered here suggests that the
artefact has more institutional texture than a policy-advocacy reading
recognises: a logic order templated centrally and reproduced across
domains; a set of legitimacy strategies that work pre-discursively and
through designed objects rather than through argument; and a
developmental-state instrument narrated in liberal-market vocabulary.
The texture is not a flaw to be corrected. It is the shape an EU
mission-oriented instrument takes when it is built inside an EU
procurement architecture and addressed to an audience for whom
industrial policy still has to be made acceptable. What this thesis
contributes is the vocabulary to see the texture, and a set of
points in the text at which that vocabulary applies.
